\section{Results}

The Roney deposit contributed $100$ left-atrial meshes and all $100$ were
analysed. The University of Washington (UW)/Boyle deposit contained $164$
meshes, being $82$ patients imaged before and after ablation; the $82$
pre-ablation meshes were analysed, since post-ablation geometry encodes the
delivered treatment. No subject was excluded for any other reason, so the
analysed real-anatomy cohort was $100+82=182$ with no attrition. The
electrophysiology simulator returned a reentry-inducible verdict in $34$ of
$100$ Roney subjects ($34.0\%$) and in $8$ of $82$ UW subjects ($9.8\%$), $42$
of $182$ combined; the UW rate sits far below the moderate-burst band expected
for a cohort of atrial fibrillation patients and is flagged as anomalous. Up to
$100{,}000$ synthetic excitable networks supplied the scale arm. Table~1 gives
cohort composition, event counts, baseline discrimination, and the incremental
results.

The pre-registered primary endpoint, a grouped $\dAUC\geq0.05$ with DeLong
$p<0.05$ when \SFI{} was added to the six published competitor features, was
\textbf{not met} in any cohort or by either classifier. On the combined real
cohort the competitor baseline reached an out-of-fold AUC of $0.884$ under
logistic regression and $0.851$ under gradient boosting; adding the \SFI{} block
changed discrimination by $\dAUC=-0.007$ ($95\%$ CI $-0.042$ to $+0.029$;
$p=0.70$) and $+0.012$ ($95\%$ CI $-0.017$ to $+0.044$; DeLong $p=0.449$),
respectively. The gradient-boosting point estimate is roughly a quarter of the
pre-registered threshold and its interval excludes that threshold as well as
containing zero, so the combined result is a measured null rather than a
shortfall of power. Cohort-specific increments were $-0.021$ and $+0.009$ on
Roney and $+0.000$ and $+0.035$ on UW. The UW arm rests on eight positives and
its intervals ($[-0.073,+0.081]$ and $[-0.095,+0.203]$) reach past the
effect-size gate, so that arm is uninformative rather than null (Figure~1). The
pooled row is a single ROC across two cohorts of unequal prevalence, so it
carries between-cohort as well as within-cohort separability and does not
replace the cohort-specific estimates.

The null held at scale. On $2{,}000$ FitzHugh--Nagumo networks the increment was
$\leq+0.003$ with $p\geq0.12$, and on $500$ Kuramoto phase-oscillator networks
$\leq+0.0036$ with $p\geq0.38$. Extending the FitzHugh--Nagumo cohort to
$100{,}000$ networks resolved the effect statistically without enlarging it: the
subspace variant added $+0.0028$ at $p<10^{-4}$ once $N\geq10^{4}$. Detecting an
increment of that size at $80\%$ power would require approximately $4{,}150$
cases, and the single-vector variant's $+0.0007$ approximately $23{,}000$. That
characterisation was not transported to the real cohorts, a separate regime.

Because a zero increment is equally consistent with redundancy, with an
uninformative feature, and with one the classifier cannot exploit, redundancy was
measured directly and without reference to the label. The leading canonical
correlation between the nine \SFI{} columns and the six competitors was $0.996$,
against $0.289$ for a permuted-\SFI{} control, so the two blocks are almost
collinear in their dominant direction. The per-column picture was less uniform:
the total and region-mean summaries were reconstructible from the competitors at
$R^{2}=0.970$, four of the nine columns reached $R^{2}\geq0.60$, and the
remaining five --- the region maximum, the region standard deviation, the top
single region, and both edge-level columns --- retained most of their variance at
$R^{2}<0.25$, well above the permuted control. Redundancy therefore accounts for
part of the null but not all of it: for some columns the competitors already
carry the information, while for others unique variance exists and does not
predict the label. Decisively, the \emph{exact} per-region $\Delta\lam_2$
variant, which carries no first-order approximation error, also added nothing
over the competitor set, so the null is not a failure of the approximation.

Estimator validity was assessed without reference to any label. The
single-vector first-order form is trustworthy only while
$\rhoval=\lVert\Delta L\rVert/(\lam_3-\lam_2)$ is $O(1)$; empirically its
relative error is $1\%$ at $\rhoval=0.3$ and crosses $10\%$ at
$\rhoval^\ast\approx3$. Across the $12$ scanned real atrial meshes the median was
$\rhoval\approx2422$ (interquartile range $[1082,7384]$; full range
$[817,16554]$), so every mesh exceeded $\rhoval^\ast$ and even the smallest did so
by a factor of $272$. The spectral gap in the denominator collapses lawfully with
resolution: measured on the weighted conduction Laplacian the biomarker itself
uses, the atrial surface fitted $N^{-1.00}$, within $0.5\%$ of the Weyl
2-manifold exponent of $-1.0$, so meshing the same anatomy more finely degrades
the estimator further rather than improving it (Figure~2). This establishes that
the single-vector estimator is numerically invalid on real tissue and that
$\phii_2$ is ill-conditioned; it does not by itself establish
non-predictiveness.

The only endpoint scored against patient outcomes was two-year post-ablation
arrhythmia recurrence in the $82$ UW patients, $48$ of whom recurred, with every
patient followed for the full two years (Table~2). At $n=82$ with $48$ events,
power at $\dAUC=0.05$ was $0.24$, so the analysis plan --- frozen and tagged
before any feature was placed beside the outcome column --- fixed the minimum
detectable effect at $\dAUC\approx0.10$--$0.15$, specified a fixed-sequence
stopping rule, and pre-committed that a materially degraded baseline would be
reported as confounding rather than as a win for \SFI{}. Under logistic
regression the competitor baseline gave an out-of-fold AUC of $0.250$, adding
\SFI{} gave $0.447$, and the increment was $\dAUC=+0.197$ ($95\%$ CI $+0.054$ to
$+0.334$; DeLong $p=0.0061$), clearing both arithmetic thresholds. The
pre-specified reading is that the increment is an artefact. A baseline AUC of
$0.250$ is anti-predictive rather than weak, and $0.447$ remains below chance.
All $15$ features, competitor and \SFI{} alike, were univariately at chance
($0.435$--$0.547$); the gradient-boosting arm, whose baseline sat at $0.516$,
gave $\dAUC=+0.021$ ($p=0.738$); and replacing the nine \SFI{} columns with nine
columns of standard Gaussian noise and re-running the identical pipeline
reproduced the increment, the noise gaining $+0.137$ on average across $20$ draws
with $6$ of $20$ matching or exceeding $+0.197$ (empirical $p=0.333$). Fold event
rates ranged from $25\%$ to $81\%$, and mean predicted probability on a held-out
fold tracked that fold's training prevalence at $r=+0.995$, so pooling
out-of-fold probabilities inverted the ranking; scoring within folds and
averaging returned the baseline to $0.345$. The clinical primary endpoint was
\textbf{not met}, the observed increment is reported as confounded by baseline
degradation, and under the stopping rule the ablation-induced $\Delta$\SFI{} test
was neither performed nor computed.

Removing any single pre-registered guard manufactured a specific false positive
(Figure~3). On a controlled replicated cohort, replacing shape-family grouping
with naive random cross-validation inflated AUC from $0.640$ to $0.756$, a gain
of $+0.117$. Removing the effect-size gate admitted the $p$-value trap: at
$N=55{,}000$ networks the subspace \SFI{} reached a paired DeLong
$p\approx2\times10^{-24}$ over the competitor set while its $\dAUC$ of $+0.003$
would require approximately $4{,}150$ cases to detect at $80\%$ power.
